%!TEX TS-program = xelatex
%!TEX encoding = UTF-8 Unicode
\documentclass[11pt,a4paper]{awesome-cv}
\usepackage[hidelinks]{hyperref}

% --- ATS / PDF metadata ---
\hypersetup{
  pdftitle={Juan Manuel Young Hoyos — Cover Letter — Mobile Engineer},
  pdfauthor={Juan Manuel Young Hoyos},
  pdfsubject={Cover Letter — Mobile Engineer — Sezzle},
  pdfkeywords={React Native, TypeScript, mobile, App Store, Google Play, automated tests, Claude, LLM, Golang, Python, fintech, payments, Sezzle}
}

\geometry{left=1.4cm, top=.8cm, right=1.4cm, bottom=1.8cm, footskip=.5cm}

% Emit ActualText so PDF text extraction yields ASCII hyphens (ATS keyword matching)
\XeTeXgenerateactualtext=1

\colorlet{awesome}{awesome-red}
\setbool{acvSectionColorHighlight}{true}
\renewcommand{\acvHeaderSocialSep}{\quad\textbar\quad}

\definecolor{darktext}{HTML}{1A1A1A}
\colorlet{text}{darktext}
\renewcommand*{\lettertextstyle}{\fontsize{10pt}{1.45em}\bodyfontlight\upshape\color{darktext}}
\renewcommand*{\lettersectionstyle}[1]{{\fontsize{14pt}{1em}\bodyfont\bfseries\color{darktext}#1}}

% --- Header info ---
\name{Juan Manuel}{Young Hoyos}
\position{Mobile Engineer — React Native · TypeScript · AI-Assisted Development}
\address{Medellín, Colombia · 100\% Remote · English B2+}
\mobile{+57 1234567890}
\email{juanmanuel12.13jmyh81@gmail.com}
\homepage{jmyounghoyos.com}
\github{Youngermaster}
\linkedin{juan-manuel-young-hoyos}

% --- Letter meta ---
\recipient
  {Engineering — Hiring Team}
  {Sezzle}
\letterdate{\today}
\lettertitle{Application for Mobile Engineer (Colombia)}
\letteropening{Dear Sezzle Hiring Team,}
\letterclosing{Sincerely,}
\letterenclosure[Attached]{Resume — Juan Manuel Young Hoyos.pdf}

\begin{document}
\makecvheader[R]
\makelettertitle

\begin{cvletter}

I am applying for the \textbf{Mobile Engineer} role at Sezzle. I am based in Medellín, I work fully remotely with US teams already, and I have spent the last few years doing precisely what this role describes: building features in a high-traffic, cross-platform \textbf{TypeScript React Native} app.

At \textbf{Publicis Global Delivery} I built production React Native features for \textbf{Wegmans Food Markets}, one of the largest US grocery retailers. Modular TypeScript components with \textbf{Zustand} state management and Algolia search lifted \textbf{in-app engagement 30\%}; comprehensive \textbf{Jest unit tests} and cross-device verification with BrowserStack cut bug incidence \textbf{25\%}; and helping build a \textbf{CI/CD pipeline for Android APK generation} cut release-cycle time \textbf{50\%}. Through Grisú I have also designed and \textbf{launched} my own React Native apps end to end, so I know the release path, not just the feature work.

Your requirement around \textbf{Claude and LLM tooling} is one I meet genuinely rather than aspirationally --- I use \textbf{Claude Code} daily to build, refactor, generate tests, research and document, and I review and validate everything it produces. Speed matters, but defects should not get sent down the line, so AI raises my throughput without lowering the bar.

I also fit the wider stack. I have written backend services in \textbf{Golang} (Go Fiber) and \textbf{Python} (FastAPI), shipped \textbf{React} on the web, deployed with Docker on AWS and Azure, and used \textbf{GitLab} CI in my most recent role --- so I can contribute across a full-stack team rather than only inside the app. On fintech specifically: my team won \textbf{1st place among 52+ teams from 16 universities} at the CivicaPay hackathon with a working mobile payments prototype, which is where my interest in payments started.

The parts of your culture description that resonate most are having relentlessly high standards and disagreeing then committing. I would rather raise a concern early and be wrong than stay quiet and be right later. Thank you for your consideration.

\end{cvletter}

\makeletterclosing
\end{document}
